% Appendix E: Carrier Forcing Proofs (Chapters 7--8)
\label{sec:appendix-carrier}

% ── Chapter 7: ch7_carrier_forcing.tex ──

\begin{proof}[Proof of Theorem~\ref{thm:persistence-mode-separation}]
\label{proof:thm:persistence-mode-separation}
Local-state persistence and propagative persistence are distinct because the
former preserves phase-compatible reversible scalar transport, while the latter
preserves reversible second-order transport with no preferred direction. These
are different compatibility signatures. Comparative persistence is distinct
from both because its physically visible datum is a comparison class modulo
local relabeling, governed by exactness and local gauge covariance rather than
by scalar transport alone. The causal-structural lane is distinct from
comparative persistence because it is not the same visible datum at all: it is
the symmetric-response lift of the primitive \(3\)-channel, with
divergence-compatible response and tensor exchange source, not a one-form
comparison class with exactness constraint. In each case, identifying two
roles would remove one of the defining compatibility signatures or replace one
admissible source profile by another. Therefore the three primitive channels
and the causal-structural lift are pairwise distinct in the required sense.
\end{proof}

\begin{proof}[Proof of Theorem~\ref{thm:resolution_integrated_balance_closure_not_intrinsic_mode}]
\label{proof:thm:resolution_integrated_balance_closure_not_intrinsic_mode}
The intrinsic roles of Definition~\ref{def:five-persistence-modes} are defined
either by direct compatibility signatures visible on irreducible subcarriers of
the selected cut itself, or by the canonical symmetric-response lift of the
primitive \(3\)-channel: phase-compatible scalar transport, reversible
second-order propagation, exactness with local gauge covariance, and the
geometric symmetric-response lift. A hydrodynamic law obtained from a
resolution-integrated balance hierarchy is different in kind. Its exact visible
content begins with the balance hierarchy on collision-invariant observables,
while its recognizable Navier--Stokes-type form appears only after unresolved
exchange has been compressed into an effective endpoint response. Therefore it
is not an intrinsic visible law on the same footing as the three primitive
channels or the geometric lift, but an effective closure built from integrated
unresolved exchange.
\end{proof}

\begin{proof}[Proof of Theorem~\ref{thm:canonical-sector-generation-reduction}]
\label{proof:thm:canonical-sector-generation-reduction}
Hypothesis~\emph{(i)} identifies the selected visible effective law as the
exact reduced law after residual return elimination, so all observable forcing
already factors through returned residual on the selected cut. Hypothesis
~\emph{(ii)} gives one common selected observed bundle carrying the four
intrinsic flagship sectors together with the hydrodynamic closure lane.
Hypothesis~\emph{(iii)} then states that these intrinsic sectors are not
assigned by a primitive dictionary but realized as the canonical minimal
compatible realizations of the three primitive persistence channels together
with the canonical symmetric-response lift of the primitive \(3\)-channel.
Hypothesis~\emph{(iv)} then records the kinetic lane in its proper place: not
as a further primitive channel or intrinsic lift, but as the
resolution-integrated effective balance closure carried by the same selected
bundle. Once those realizations are fixed, one identifies the unique canonical
representative of the distinguished minimal compatible closure class on each
intrinsic carrier together with the canonical effective normal form on the
balance hierarchy, which is exactly the data requested by
Definition~\ref{def:carrier-level-classification}.
\end{proof}

\begin{proof}[Proof of Theorem~\ref{thm:existence-of-five-persistence-modes}]
\label{proof:thm:existence-of-five-persistence-modes}
By Theorem~\ref{thm:canonical-sector-generation-reduction}, the four intrinsic
flagship carriers on the selected observed bundle are exactly the canonical
realizations of the persistence modes listed in
Definition~\ref{def:canonical-sector-generation}. Those realization clauses are
explicit: the phase carrier realizes local-state persistence, the wave carrier
realizes propagative persistence, the gauge carrier realizes comparative
persistence, and the geometric carrier realizes causal-structural persistence
as the canonical symmetric-response lift of the primitive \(3\)-channel. Hence
every one of the three primitive channels together with the geometric lift is
realized on the selected observed bundle.
\end{proof}

\begin{proof}[Proof of Theorem~\ref{thm:flagship-mode-classification}]
\label{proof:thm:flagship-mode-classification}
By Theorem~\ref{thm:canonical-sector-generation-reduction}, the intrinsic
flagship lanes are precisely the canonical realizations of the structural
roles listed in Definition~\ref{def:canonical-sector-generation} on the
selected observed bundle. Thus each such lane belongs to at least one of those
roles. By Theorem~\ref{thm:persistence-mode-separation}, the three primitive
channels and the geometric lift are pairwise distinct in role, so no intrinsic
flagship lane can belong to more than one without destroying its defining
compatibility signature or source-admissibility profile. Therefore each
intrinsic flagship visible law belongs to exactly one of those structural
roles.
\end{proof}

\begin{proof}[Proof of Corollary~\ref{cor:flagship-family-exhausts-five-persistence-modes}]
\label{proof:cor:flagship-family-exhausts-five-persistence-modes}
Item~\emph{(i)} is exactly
Theorem~\ref{thm:flagship-mode-classification}. For item~\emph{(ii)},
Theorem~\ref{thm:existence-of-five-persistence-modes} gives at least one
intrinsic flagship realization of each primitive channel together with the
geometric lift. Suppose one of the three primitive channels were realized by
two distinct intrinsic flagship carrier classes. By
Definition~\ref{def:canonical-sector-generation}, those distinct carrier
classes come with distinct compatibility signatures on the selected observed
bundle. One of those carrier classes would then belong both to its own
designated mode and to the repeated mode, contradicting
Theorem~\ref{thm:flagship-mode-classification}. Hence each primitive channel is
realized by exactly one intrinsic flagship carrier class. The geometric clause
of item~\emph{(ii)} is exactly the fourth clause of
Definition~\ref{def:canonical-sector-generation}. For item~\emph{(iii)}: if a
fourth primitive channel or an additional intrinsic lift occurred in the
intrinsic flagship family, some intrinsic flagship carrier class would realize
a role outside the three primitive channels together with the geometric lift
listed in Definition~\ref{def:five-persistence-modes},
contradicting item~\emph{(i)}.
\end{proof}

\begin{proof}[Proof of Theorem~\ref{thm:common-selected-bundle-assembly}]
\label{proof:thm:common-selected-bundle-assembly}
These four intrinsic carriers together with the hydrodynamic balance carrier
are taken on the same physically realized selected cut, so by
Definition~\ref{def:characteristic-observer-datum} they are visible sectors of
the one selected bundle \(\mathcal X_h=P_h\cH^\bullet\). By
Definition~\ref{def:selected-visible-effective-law}, the same exact effective
operator \(\mathcal L_{\mathrm{vis},h}(z)\) and the same inherited endpoint
closure mechanism are evaluated on that cut before any carrier-specific
restriction is made. The only remaining carrier-dependent step is therefore the
classification of the surviving endpoint operator on the chosen visible
carrier, followed by the final carrier-level physical identification. This is
exactly the common grammar of Remark~\ref{rem:common-part-iv-grammar}.
\end{proof}

\begin{proof}[Proof of Theorem~\ref{thm:common-endpoint-rigidity-reduction}]
\label{proof:thm:common-endpoint-rigidity-reduction}
By
Theorem~\ref{thm:least-motion-observer-equivalence-class-is-admissible-selected-class}(iv),
together with Theorem~\ref{thm:selection-interface} and
Corollary~\ref{cor:selection-to-forcing}, the selected observed law on
\(\mathcal X_h^{\mathrm{vis}}\) is unique modulo horizon-preserving
equivalence at each relevant horizon. The scalarization surface built into the
theorem statement, together with
Corollary~\ref{cor:realized-self-adjoint-scalarization} and
Corollary~\ref{cor:realized-isotropic-symmetric-response-scalarization},
reduces every admissible selected closure operator to scalar action on the
intrinsic symmetry channels carried by \(\mathcal X_h^{\mathrm{vis}}\). By
Definition~\ref{def:selected-visible-effective-law}, the selected endpoint
closure class on that same carrier is already the carrier-side exact closure
class obtained from the one exact visible law on the selected cut by
characteristic boundary reduction. Hence that class already has a unique exact
carrier-side representative and is unique modulo the continuum equivalence used
later in the carrier-level endpoint identifications. Combined with the
channelwise scalarization above, the class is encoded entirely by its scalar
data on the intrinsic symmetry channels. The only remaining carrier-level
physical task is therefore to identify the canonical representative of that
already forced class.
\end{proof}

% ── Chapter 8: flagship_quantum.tex (Phase / Schr\"odinger) ──

\begin{proof}[Proof of Theorem~\ref{thm:least-motion-phase-observer}]
\label{proof:thm:least-motion-phase-observer}
Hypothesis \emph{(i)} gives faithfulness, closure-admissibility, and promotion
stability of the phase carrier. Hypothesis \emph{(ii)} says that no proper
visible subcarrier still supports the same forced phase, transport, and scalar
closure data. Hypothesis \emph{(iii)} says that any competing characteristic
observer datum with those properties is horizon-preserving equivalent to the
datum generated by \(\mathcal X_{\mathrm{ph}}\). Therefore that datum
satisfies the clauses of
Definition~\ref{def:least-motion-characteristic-observer} and is unique up to
horizon-preserving equivalence. Hence \(\mathcal X_{\mathrm{ph}}\) is the
unique least-motion characteristic observer for the realized family.
\end{proof}

\begin{proof}[Proof of Theorem~\ref{thm:induced-phase-hamiltonian-law}]
\label{proof:thm:induced-phase-hamiltonian-law}
By Theorem~\ref{thm:least-motion-phase-observer}, the visible law is carried by
the selected phase carrier \(\mathcal X_{\mathrm{ph}}\). For each discrete
level of the realized family,
Corollary~\ref{cor:intrinsic-package-implies-governing-datum-rigid} places the
selected channels in a governing-datum rigid realized class. Therefore
Theorem~\ref{cor:realized-local-internal-datum-forcing} yields, on every
relevant selected channel, a forced phase-carrier class, a unique nonnegative
scalar quadratic recursion datum, and a scalar closure class.

By the phase-Hamiltonian realization hypotheses, those selected channel data
are organized into a refinement-compatible family on
\(C^1(J,\mathcal W_0)\) with a refinement-coherent reconstruction datum whose
exact core realization is bounded. The transport side therefore carries the
self-adjoint second-order operator \(\mathcal T_{\chi}\), while the scalar
closure side carries the self-adjoint operator \(V_{\chi}\) on the same
selected cut. The phase side inherits a representative \(J_{\chi}\) of the
forced phase-carrier class on that cut. The unique scalar coefficient induced
by the selected transport class is denoted \(\kappa_{\chi}\ge 0\). All
higher-order unresolved endpoint return is absorbed into the remainder
\(\mathsf R_{\mathrm{ph}}\), which by
Theorem~\ref{thm:temporal-realization-forbids-primitive-internal-creation}
records the admissible exchange trace of unresolved endpoint return. This
yields the displayed reversible phase law.
\end{proof}

\begin{proof}[Proof of Theorem~\ref{thm:phase-generation-and-energy-orthogonality-force-complex-generator}]
\label{proof:thm:phase-generation-and-energy-orthogonality-force-complex-generator}
By Corollary~\ref{cor:realized-phase-generation-forcing}, every relevant
selected phase channel carries a forced phase-carrier class modulo
phase-equivalence. Passing to the refinement-compatible selected carrier
\(\mathcal X_{\mathrm{ph}}\) therefore yields a unique weak phase-law class on
the selected cut, with residual freedom only by channelwise phase-equivalent
sign changes.

By
Theorem~\ref{thm:canonical-selected-schur-bridge-forces-exact-selected-cut-energy-split},
the canonical selected Schur bridge supplies the exact selected-cut energy
split and hence the canonical pairing on \(\mathcal X_{\mathrm{ph}}\). The
weak phase law of Theorem~\ref{thm:induced-phase-hamiltonian-law} already fixes
the selected transport coefficient \(\kappa_{\chi}\ge 0\), the transport class,
and the scalar closure class. Replacing the phase representative by a
nontrivial phase-equivalent sign change would reverse the phase side of that
same weak law on the affected selected channel while leaving the transport and
closure data fixed, so it would no longer represent the same weak law on the
canonical pairing. Hence exactly one member of the phase-equivalence class is
pairing-compatible with the selected weak law.

Therefore the forced phase-carrier class on the selected cut has a unique
pairing-compatible strong representative. That representative is, by
construction, the canonical complex structure carried by the selected phase
cut. Writing that canonical complex structure in standard complex notation
gives the stated complex phase generator.
\end{proof}

\begin{proof}[Proof of Theorem~\ref{thm:phase-local-endpoint-collapse}]
\label{proof:thm:phase-local-endpoint-collapse}
By
Theorem~\ref{thm:phase-generation-and-energy-orthogonality-force-complex-generator},
the forced phase-carrier class on the selected cut already has a unique
pairing-compatible strong representative, written in standard complex notation
as \(J_{\chi}=i\).

For the transport side, item \emph{(i)} gives a finite local closure surface
of scalar second-order candidates, while item \emph{(ii)} supplies the
no-preferred-direction quadratic rigidity on each relevant selected channel.
By Corollary~\ref{cor:realized-no-preferred-direction-quadratic-transport}, the
quadratic transport datum is already reduced channelwise to a single scalar
coefficient. The completion-and-selection algebra of
Section~\ref{sec:constraint-compiler} therefore collapses the admissible local
pairing-symmetric second-order scalar class to its unique isotropic local
representative. By item \emph{(iii)} and
Theorem~\ref{thm:refinement-coherent-core-realization}, the reconstructed
transport family carries a unique exact core representative; denote that
positive
Laplace-type representative by \(\Delta_{\chi}\).

For the closure side, item \emph{(iv)} says that the selected scalar closure
family has local order \(0\). Hence at each site of the local scalar carrier,
the value of the closure operator depends only on the value of the field at
that same site. Linearity therefore forces pointwise multiplication by a scalar
profile, and self-adjointness forces that profile to be real. By item
\emph{(v)} and Theorem~\ref{thm:refinement-coherent-core-realization}, the
corresponding reconstructed closure family carries a unique exact core
representative; write it as the real multiplication operator \(V\).

Substituting the forced representatives \(J_{\chi}=i\),
\(\mathcal T_{\chi}=\Delta_{\chi}\), and the scalar closure operator \(V\) into
Theorem~\ref{thm:induced-phase-hamiltonian-law} gives the displayed identity.
The lossless statement is immediate.
\end{proof}

\begin{proof}[Proof of Corollary~\ref{cor:phase-schur-remainder-identity}]
\label{proof:cor:phase-schur-remainder-identity}
By Theorem~\ref{thm:canonical-selected-schur-realization}, the only hidden-side
visible forcing on the selected cut is the returned residual, and that
returned residual is exactly the Schur correction \(\CCSigma{h}{z}\). By
Theorem~\ref{thm:phase-local-endpoint-collapse}, the selected phase endpoint
class has a unique canonical representative on the selected carrier. Therefore
the common Schur correction has a unique phase-side image on that
representative; denote it by \(\Sigma_{h,\mathrm{ph}}(z)\). By
Theorem~\ref{thm:returned-remainders-form-predictive-correction-surface}, the
phase remainder is exactly that returned correction on the phase lane, so
\(\mathsf R_{\mathrm{ph}}=\Sigma_{h,\mathrm{ph}}(z)\). Substituting this
identity into Theorem~\ref{thm:phase-local-endpoint-collapse} gives the
displayed equation.
\end{proof}

\begin{proof}[Proof of Theorem~\ref{thm:hidden-channel-phase-remainder-and-spectral-splitting}]
\label{proof:thm:hidden-channel-phase-remainder-and-spectral-splitting}
Substituting the common plane-wave ansatz of item \emph{(iv)} into the visible
and hidden phase laws gives
\[
\omega\Psi
=
E_{\mathrm{vis}}(k)\Psi+\gamma_{\chi}H,
\]
\[
\omega H
=
E_{\mathrm{hid}}(k)H+\gamma_{\chi}\Psi.
\]
If \(E_{\mathrm{hid}}(k)\neq \omega\), the second identity gives
\[
H
=
\frac{\gamma_{\chi}}{E_{\mathrm{hid}}(k)-\omega}\Psi.
\]
Substituting this into the visible equation yields
\[
\omega\Psi
=
E_{\mathrm{vis}}(k)\Psi
+
\frac{\gamma_{\chi}^2}{E_{\mathrm{hid}}(k)-\omega}\Psi.
\]
Cancelling the nonzero visible amplitude \(\Psi\) gives
\[
\omega
=
E_{\mathrm{vis}}(k)
+
\frac{\gamma_{\chi}^2}{E_{\mathrm{hid}}(k)-\omega}.
\]
This is exactly the displayed phase law with
\[
\sigma_{\chi}^{\mathrm{ph,hid}}(k,\omega)
=
\frac{\gamma_{\chi}^2}{E_{\mathrm{hid}}(k)-\omega}.
\]
Multiplying by \(E_{\mathrm{hid}}(k)-\omega\) gives
\[
(\omega-E_{\mathrm{vis}}(k))
(\omega-E_{\mathrm{hid}}(k))
=
\gamma_{\chi}^2.
\]
Solving this quadratic equation in \(\omega\) yields the two displayed branch
functions \(\omega_{\pm}(k)\).
\end{proof}

\begin{proof}[Proof of Corollary~\ref{cor:exact-off-resonant-phase-line-shifts}]
\label{proof:cor:exact-off-resonant-phase-line-shifts}
By
Theorem~\ref{thm:hidden-channel-phase-remainder-and-spectral-splitting},
\[
\omega_{\pm}(k)
=
\frac{
E_{\mathrm{vis}}(k)+E_{\mathrm{hid}}(k)
\pm
\sqrt{
\bigl(\delta_{\chi}^{\mathrm{ph}}(k)\bigr)^2+4\gamma_{\chi}^2
}
}{2}.
\]
Since \(E_{\mathrm{hid}}(k)=E_{\mathrm{vis}}(k)+\delta_{\chi}^{\mathrm{ph}}(k)\),
\[
\omega_-(k)
=
E_{\mathrm{vis}}(k)
+
\frac{
\delta_{\chi}^{\mathrm{ph}}(k)
-
\sqrt{
\bigl(\delta_{\chi}^{\mathrm{ph}}(k)\bigr)^2+4\gamma_{\chi}^2
}
}{2}.
\]
Rationalizing the numerator gives the displayed formula for \(\omega_-(k)\),
and the formula for \(\omega_+(k)\) is obtained similarly. Because
\(\delta_{\chi}^{\mathrm{ph}}(k)>0\), the visible-like shift is negative and
the hidden-like shift is positive.
\end{proof}

\begin{proof}[Proof of Corollary~\ref{cor:exact-avoided-crossing-coupled-phase-branches}]
\label{proof:cor:exact-avoided-crossing-coupled-phase-branches}
At a selected mode \(k_\ast\) with
\(E_{\mathrm{vis}}(k_\ast)=E_{\mathrm{hid}}(k_\ast)\), the branch formula of
Theorem~\ref{thm:hidden-channel-phase-remainder-and-spectral-splitting} becomes
\[
\omega_{\pm}(k_\ast)
=
\frac{2E_\ast\pm\sqrt{4\gamma_{\chi}^2}}{2}
=
E_\ast\pm\gamma_{\chi}.
\]
Subtracting the two identities gives the exact gap.
\end{proof}

\begin{proof}[Proof of Corollary~\ref{cor:low-energy-renormalized-phase-offset-and-transport-coefficient}]
\label{proof:cor:low-energy-renormalized-phase-offset-and-transport-coefficient}
For small \(k\),
\[
\sin^2\!\left(\frac{k\ell_{\chi}}{2}\right)
=
\frac{k^2\ell_{\chi}^2}{4}+O(k^4),
\]
so
\[
E_{\mathrm{vis}}(k)
=
U_{\chi}+\kappa_{\chi}k^2+O(k^4),
\qquad
E_{\mathrm{hid}}(k)
=
\widetilde U_{\chi}+\widetilde\kappa_{\chi}k^2+O(k^4).
\]
Hence
\[
\delta_{\chi}^{\mathrm{ph}}(k)
=
\Delta_0^{\mathrm{ph}}+\Delta_1^{\mathrm{ph}}k^2+O(k^4).
\]
By
Theorem~\ref{thm:hidden-channel-phase-remainder-and-spectral-splitting},
\[
\omega_-(k)
=
\frac{
E_{\mathrm{vis}}(k)+E_{\mathrm{hid}}(k)
-
\sqrt{
\bigl(\delta_{\chi}^{\mathrm{ph}}(k)\bigr)^2+4\gamma_{\chi}^2
}
}{2}.
\]
Expanding the square root at \(k=0\) gives
\[
\sqrt{
\bigl(\delta_{\chi}^{\mathrm{ph}}(k)\bigr)^2+4\gamma_{\chi}^2
}
=
\sqrt{
\bigl(\Delta_0^{\mathrm{ph}}\bigr)^2+4\gamma_{\chi}^2
}
+
\frac{
\Delta_0^{\mathrm{ph}}\Delta_1^{\mathrm{ph}}
}{
\sqrt{
\bigl(\Delta_0^{\mathrm{ph}}\bigr)^2+4\gamma_{\chi}^2
}
}k^2
+
O(k^4).
\]
Substituting these expansions into the branch formula and collecting the
constant and quadratic terms yields the displayed formulas for
\(U_{\chi,\mathrm{eff}}\) and \(\kappa_{\chi,\mathrm{eff}}\).
\end{proof}

% ── Chapter 8: flagship_wave.tex (Wave / Klein--Gordon) ──

\begin{proof}[Proof of Theorem~\ref{thm:least-motion-wave-observer}]
\label{proof:thm:least-motion-wave-observer}
Hypothesis \emph{(i)} gives faithfulness, closure-admissibility, and promotion
stability of the wave carrier. Hypothesis \emph{(ii)} says that no proper
visible subcarrier still supports the same scalar second-order transport and
scalar endpoint closure data. Hypothesis \emph{(iii)} says that any competing
characteristic observer datum with those properties is horizon-preserving
equivalent to the datum generated by \(\mathcal X_{\mathrm{wav}}\).
Therefore that datum satisfies the clauses of
Definition~\ref{def:least-motion-characteristic-observer} and is unique up to
horizon-preserving equivalence. Hence \(\mathcal X_{\mathrm{wav}}\) is the
unique least-motion characteristic observer for the realized family.
\end{proof}

\begin{proof}[Proof of Theorem~\ref{thm:induced-wave-law}]
\label{proof:thm:induced-wave-law}
By Theorem~\ref{thm:least-motion-wave-observer}, the visible law is carried by
the selected wave carrier \(\mathcal X_{\mathrm{wav}}\). For each discrete
level of the realized family,
Corollary~\ref{cor:intrinsic-package-implies-governing-datum-rigid} places the
selected channels in a governing-datum rigid realized class. Therefore
Theorem~\ref{cor:realized-local-internal-datum-forcing} yields, on every
relevant selected channel, a forced scalar second-order transport datum and a
scalar endpoint closure class.

By the reversible wave realization hypotheses, those selected channel data are
organized into a refinement-compatible family on \(C^2(J,\mathcal W_0)\) with
a refinement-coherent reconstruction datum whose exact core realization is
bounded. The transport side therefore carries the self-adjoint nonnegative
second-order operator \(\mathcal T_{\chi}\), while the scalar closure side
carries the self-adjoint operator \(M_{\chi}\) on the same selected cut. The
unique scalar coefficient induced by the selected transport class is denoted
\(c_{\chi}>0\). All higher-order unresolved endpoint return is absorbed into
the remainder \(\mathsf R_{\mathrm{wav}}\), which by
Theorem~\ref{thm:temporal-realization-forbids-primitive-internal-creation}
records the admissible exchange trace of unresolved endpoint return. This
yields the displayed reversible wave law.
\end{proof}

\begin{proof}[Proof of Theorem~\ref{thm:wave-local-endpoint-collapse}]
\label{proof:thm:wave-local-endpoint-collapse}
For the transport side, item \emph{(i)} gives a finite local closure surface
of scalar second-order candidates, while item \emph{(ii)} supplies the
no-preferred-direction quadratic rigidity on each relevant selected wave
channel. By Corollary~\ref{cor:realized-no-preferred-direction-quadratic-transport},
the quadratic transport datum is already reduced channelwise to a single scalar
coefficient. The completion-and-selection algebra of
Section~\ref{sec:constraint-compiler} therefore collapses the admissible local
pairing-symmetric second-order scalar class to its unique isotropic local
representative. By item \emph{(iii)} and
Theorem~\ref{thm:refinement-coherent-core-realization}, the reconstructed
transport family carries a unique exact core representative; denote the corresponding
positive Laplace-type representative by \(\Delta_{\chi}\). In the sign
convention of Theorem~\ref{thm:induced-wave-law}, this means
\[
\mathcal T_{\chi}=-\Delta_{\chi}.
\]

For the closure side, item \emph{(iv)} says that the selected scalar closure
family has local order \(0\). Hence at each site of the local scalar carrier,
the value of the closure operator depends only on the value of the field at
that same site. Linearity therefore forces pointwise multiplication by a scalar
profile, and self-adjointness forces that profile to be real. By item
\emph{(v)} and Theorem~\ref{thm:refinement-coherent-core-realization}, the
corresponding reconstructed closure family carries a unique exact core
representative; write it as the real multiplication operator \(W_{\chi}\).

Substituting \(\mathcal T_{\chi}=-\Delta_{\chi}\) and \(M_{\chi}=W_{\chi}\)
into Theorem~\ref{thm:induced-wave-law} yields the displayed identity. The
lossless statement is immediate.
\end{proof}

\begin{proof}[Proof of Corollary~\ref{cor:wave-schur-remainder-identity}]
\label{proof:cor:wave-schur-remainder-identity}
By Theorem~\ref{thm:canonical-selected-schur-realization}, the selected cut
carries one common returned residual field, and that returned residual is
exactly the Schur correction \(\CCSigma{h}{z}\). By
Theorem~\ref{thm:wave-local-endpoint-collapse}, the selected wave endpoint
class has a unique canonical representative on the selected carrier. Hence the
common Schur correction has a unique wave-side image there; denote it by
\(\Sigma_{h,\mathrm{wav}}(z)\). By
Theorem~\ref{thm:returned-remainders-form-predictive-correction-surface}, the
wave remainder is exactly that returned correction on the wave lane, so
\(\mathsf R_{\mathrm{wav}}=\Sigma_{h,\mathrm{wav}}(z)\). Substituting this into
Theorem~\ref{thm:wave-local-endpoint-collapse} yields the displayed equation.
\end{proof}

\begin{proof}[Proof of Theorem~\ref{thm:homogeneous-wave-scalar-profile-forces-constant-mass}]
\label{proof:thm:homogeneous-wave-scalar-profile-forces-constant-mass}
By Theorem~\ref{thm:wave-local-endpoint-collapse}, the selected scalar closure
class is already represented on the selected wave carrier by a real
multiplication operator \(W_{\chi}\) of local order \(0\).

Fix a relevant selected wave channel. By item \emph{(i)} and
Corollary~\ref{cor:locally-homogeneous-implies-symmetry-generated}, the
corresponding admissible realized class is symmetry-generated on that channel.
By Corollary~\ref{cor:symmetry-generated-implies-closure-rigid}, the selected
self-adjoint closure operator on that homogeneous sector is closure-rigid.
Because \(W_{\chi}\) is already a zeroth-order real multiplication operator and
is equivariant under the induced local symmetry action by item \emph{(ii)},
closure-rigidity forces that multiplication profile to be scalar on the
homogeneous selected-wave sector. Denote the resulting real scalar by
\(m_{\chi}^2\).

Item \emph{(iii)} says that after restricting to the selected wave carrier
there is no residual site-dependent scalar defect beyond this homogeneous
sector. Therefore the full forced scalar profile on the selected wave carrier
is exactly constant multiplication by \(m_{\chi}^2\):
\[
W_{\chi}=m_{\chi}^2 I.
\]
Item \emph{(iv)} gives \(m_{\chi}^2\ge 0\). Uniqueness is immediate, since two
constant scalar multiplication operators agreeing on the selected wave carrier
have the same scalar.

Substituting \(W_{\chi}=m_{\chi}^2 I\) into
Theorem~\ref{thm:wave-local-endpoint-collapse} yields the displayed identity.
\end{proof}

\begin{proof}[Proof of Theorem~\ref{thm:exact-wave-dispersion}]
\label{proof:thm:exact-wave-dispersion}
Under the nearest-neighbor wave specialization, the lossless massless law of
Theorem~\ref{thm:induced-wave-law} becomes
\[
\partial_t^2\phi_n
+
\frac{c_{\chi}^2}{\ell_{\chi}^2}
\bigl(2\phi_n-\phi_{n+1}-\phi_{n-1}\bigr)
=
0.
\]
Substituting the plane-wave ansatz
\(\phi_n(t)=e^{i(nk\ell_{\chi}-\omega t)}\) gives
\[
-\omega^2
+
\frac{c_{\chi}^2}{\ell_{\chi}^2}
\bigl(2-e^{ik\ell_{\chi}}-e^{-ik\ell_{\chi}}\bigr)
=
0.
\]
Using
\[
2-e^{ik\ell_{\chi}}-e^{-ik\ell_{\chi}}
=
4\sin^2\!\left(\frac{k\ell_{\chi}}{2}\right),
\]
we obtain
\[
\omega(k)^2
=
\frac{4c_{\chi}^2}{\ell_{\chi}^2}
\sin^2\!\left(\frac{k\ell_{\chi}}{2}\right).
\]
Taking the positive branch,
\[
\omega(k)
=
\frac{2c_{\chi}}{\ell_{\chi}}
\sin\!\left(\frac{k\ell_{\chi}}{2}\right),
\]
so
\[
v_g(k)
=
\frac{d\omega}{dk}
=
c_{\chi}\cos\!\left(\frac{k\ell_{\chi}}{2}\right).
\]
Since
\[
E=\hbar\omega
=
2E_{\chi}\sin\!\left(\frac{k\ell_{\chi}}{2}\right),
\]
we have
\[
\sin\!\left(\frac{k\ell_{\chi}}{2}\right)=\frac{E}{2E_{\chi}},
\]
and therefore
\[
v_g(E)
=
c_{\chi}\sqrt{1-\left(\frac{E}{2E_{\chi}}\right)^2}.
\]
\end{proof}

\begin{proof}[Proof of Corollary~\ref{cor:observable-wave-carrier-signature}]
\label{proof:cor:observable-wave-carrier-signature}
Expand the square root in Theorem~\ref{thm:exact-wave-dispersion} at
\(E/E_{\chi}=0\):
\[
\sqrt{1-\left(\frac{E}{2E_{\chi}}\right)^2}
=
1-\frac18\left(\frac{E}{E_{\chi}}\right)^2
+O\!\left(\frac{E^4}{E_{\chi}^4}\right).
\]
Subtract \(1\) and multiply by \(c_{\chi}\). The exact cutoff is the endpoint
\(E=2E_{\chi}\) of the same formula.
\end{proof}

\begin{proof}[Proof of Theorem~\ref{thm:exact-schur-corrected-wave-dispersion}]
\label{proof:thm:exact-schur-corrected-wave-dispersion}
By Theorem~\ref{thm:homogeneous-wave-scalar-profile-forces-constant-mass} and
Corollary~\ref{cor:wave-schur-remainder-identity}, every
\(\phi\in C^2(J,\mathcal W_0)\) satisfies
\[
\partial_t^2\phi
-
c_{\chi}^2\Delta_{\chi}\phi
+
m_{\chi}^2\phi
=
\Sigma_{h,\mathrm{wav}}(z).
\]
Under item \emph{(ii)} of
Definition~\ref{def:mode-compatible-schur-corrected-lattice-wave-specialization},
\[
(\Delta_{\chi}\phi)_n
=
\frac{\phi_{n+1}-2\phi_n+\phi_{n-1}}{\ell_{\chi}^2}.
\]
Substituting the plane-wave ansatz from item \emph{(iv)} and the returned
remainder profile from item \emph{(v)} therefore gives
\[
-\omega^2\phi_n
-
\frac{c_{\chi}^2}{\ell_{\chi}^2}
\bigl(e^{ik\ell_{\chi}}-2+e^{-ik\ell_{\chi}}\bigr)\phi_n
+
m_{\chi}^2\phi_n
=
\sigma_{\chi}(k;z)\phi_n.
\]
Using
\[
2-e^{ik\ell_{\chi}}-e^{-ik\ell_{\chi}}
=
4\sin^2\!\left(\frac{k\ell_{\chi}}{2}\right),
\]
and cancelling the nonzero plane-wave factor \(\phi_n\), we obtain
\[
\omega(k;z)^2
=
m_{\chi}^2
+
\frac{4c_{\chi}^2}{\ell_{\chi}^2}
\sin^2\!\left(\frac{k\ell_{\chi}}{2}\right)
-
\sigma_{\chi}(k;z).
\]

If \(k\mapsto \sigma_{\chi}(k;z)\) is differentiable and the right-hand side
is positive, then differentiating the dispersion identity gives
\[
2\omega(k;z)\,v_g(k;z)
=
\frac{2c_{\chi}^2}{\ell_{\chi}}\sin(k\ell_{\chi})
-
\partial_k\sigma_{\chi}(k;z),
\]
which yields the displayed group-velocity formula after substituting
\[
\omega(k;z)
=
\sqrt{
m_{\chi}^2
+
\frac{4c_{\chi}^2}{\ell_{\chi}^2}
\sin^2\!\left(\frac{k\ell_{\chi}}{2}\right)
-
\sigma_{\chi}(k;z)
}.
\]
The lossless and massless reductions are immediate.
\end{proof}

\begin{proof}[Proof of Theorem~\ref{thm:hidden-channel-wave-remainder-and-branch-splitting}]
\label{proof:thm:hidden-channel-wave-remainder-and-branch-splitting}
Substituting the common plane-wave ansatz of item \emph{(iv)} into the coupled
visible and hidden equations of item \emph{(iii)} gives
\[
\bigl(\omega_{\mathrm{vis}}(k)^2-\omega^2\bigr)\Phi
=
\gamma_{\chi}H,
\]
\[
\bigl(\omega_{\mathrm{hid}}(k)^2-\omega^2\bigr)H
=
\gamma_{\chi}\Phi.
\]
If \(\omega_{\mathrm{hid}}(k)^2\neq \omega^2\), the second identity gives
\[
H
=
\frac{\gamma_{\chi}}{\omega_{\mathrm{hid}}(k)^2-\omega^2}\,\Phi.
\]
Substituting this into the first identity yields
\[
\bigl(\omega_{\mathrm{vis}}(k)^2-\omega^2\bigr)\Phi
=
\frac{\gamma_{\chi}^2}{\omega_{\mathrm{hid}}(k)^2-\omega^2}\,\Phi.
\]
Cancelling the nonzero visible amplitude \(\Phi\) gives
\[
\omega^2
=
\omega_{\mathrm{vis}}(k)^2
-
\frac{\gamma_{\chi}^2}{\omega_{\mathrm{hid}}(k)^2-\omega^2}.
\]
This is exactly the displayed visible law with
\[
\sigma_{\chi}^{\mathrm{hid}}(k,\omega)
=
\frac{\gamma_{\chi}^2}{\omega_{\mathrm{hid}}(k)^2-\omega^2}.
\]
Multiplying by \(\omega_{\mathrm{hid}}(k)^2-\omega^2\) gives
\[
\bigl(\omega_{\mathrm{vis}}(k)^2-\omega^2\bigr)
\bigl(\omega_{\mathrm{hid}}(k)^2-\omega^2\bigr)
=
\gamma_{\chi}^2,
\]
equivalently
\[
(\omega^2-\omega_{\mathrm{vis}}(k)^2)
(\omega^2-\omega_{\mathrm{hid}}(k)^2)
=
\gamma_{\chi}^2.
\]
This is a quadratic equation in \(\omega^2\), and solving it gives the two
branches \(\omega_{\pm}(k)^2\) displayed above.
\end{proof}

\begin{proof}[Proof of Corollary~\ref{cor:exact-off-resonant-hidden-channel-branch-shifts}]
\label{proof:cor:exact-off-resonant-hidden-channel-branch-shifts}
By Theorem~\ref{thm:hidden-channel-wave-remainder-and-branch-splitting},
\[
\omega_{\pm}(k)^2
=
\frac{
\omega_{\mathrm{vis}}(k)^2+\omega_{\mathrm{hid}}(k)^2
\pm
\sqrt{
\delta_{\chi}(k)^2+4\gamma_{\chi}^2
}
}{2}.
\]
Since \(\omega_{\mathrm{hid}}(k)^2=\omega_{\mathrm{vis}}(k)^2+\delta_{\chi}(k)\),
\[
\omega_-(k)^2
=
\omega_{\mathrm{vis}}(k)^2
+
\frac{\delta_{\chi}(k)-\sqrt{\delta_{\chi}(k)^2+4\gamma_{\chi}^2}}{2}.
\]
Rationalizing the numerator gives
\[
\delta_{\chi}(k)-\sqrt{\delta_{\chi}(k)^2+4\gamma_{\chi}^2}
=
-\frac{4\gamma_{\chi}^2}{
\delta_{\chi}(k)+\sqrt{\delta_{\chi}(k)^2+4\gamma_{\chi}^2}
},
\]
which yields the displayed formula for \(\omega_-(k)^2\). The formula for
\(\omega_+(k)^2\) is obtained similarly from
\[
\omega_+(k)^2
=
\omega_{\mathrm{hid}}(k)^2
-
\frac{\delta_{\chi}(k)-\sqrt{\delta_{\chi}(k)^2+4\gamma_{\chi}^2}}{2}.
\]
Because \(\delta_{\chi}(k)>0\), the correction denominator is positive, so the
visible-like shift is negative and the hidden-like shift is positive.
\end{proof}

\begin{proof}[Proof of Corollary~\ref{cor:exact-avoided-crossing-coupled-wave-branches}]
\label{proof:cor:exact-avoided-crossing-coupled-wave-branches}
At a mode \(k_\ast\) with
\(\omega_{\mathrm{vis}}(k_\ast)^2=\omega_{\mathrm{hid}}(k_\ast)^2\), the
branch formula of
Theorem~\ref{thm:hidden-channel-wave-remainder-and-branch-splitting} becomes
\[
\omega_{\pm}(k_\ast)^2
=
\frac{2\omega_\ast^2\pm\sqrt{4\gamma_{\chi}^2}}{2}
=
\omega_\ast^2\pm\gamma_{\chi}.
\]
Subtracting the two identities gives the exact gap.
\end{proof}

\begin{proof}[Proof of Corollary~\ref{cor:low-energy-renormalized-mass-and-wave-speed-hidden-channel}]
\label{proof:cor:low-energy-renormalized-mass-and-wave-speed-hidden-channel}
For small \(k\),
\[
\sin^2\!\left(\frac{k\ell_{\chi}}{2}\right)
=
\frac{k^2\ell_{\chi}^2}{4}+O(k^4),
\]
so
\[
\omega_{\mathrm{vis}}(k)^2
=
m_{\chi}^2+c_{\chi}^2k^2+O(k^4),
\qquad
\omega_{\mathrm{hid}}(k)^2
=
\widetilde m_{\chi}^2+\widetilde c_{\chi}^2k^2+O(k^4).
\]
Hence
\[
\delta_{\chi}(k)
=
\Delta_0+\Delta_1 k^2+O(k^4).
\]
By
Theorem~\ref{thm:hidden-channel-wave-remainder-and-branch-splitting},
\[
\omega_-(k)^2
=
\frac{
\omega_{\mathrm{vis}}(k)^2+\omega_{\mathrm{hid}}(k)^2
-\sqrt{\delta_{\chi}(k)^2+4\gamma_{\chi}^2}
}{2}.
\]
Expanding the square root at \(k=0\) gives
\[
\sqrt{\delta_{\chi}(k)^2+4\gamma_{\chi}^2}
=
\sqrt{\Delta_0^2+4\gamma_{\chi}^2}
+
\frac{\Delta_0\Delta_1}{\sqrt{\Delta_0^2+4\gamma_{\chi}^2}}\,k^2
+
O(k^4).
\]
Substituting the expansions into the branch formula and collecting the constant
and quadratic terms yields the displayed formulas for
\(m_{\chi,\mathrm{eff}}^2\) and \(c_{\chi,\mathrm{eff}}^2\).
\end{proof}

% ── Chapter 8: flagship_gauge.tex (Gauge / Maxwell) ──

\begin{proof}[Proof of Theorem~\ref{thm:least-motion-gauge-observer}]
\label{proof:thm:least-motion-gauge-observer}
Hypothesis \emph{(ii)} gives faithfulness, closure-admissibility, and
promotion stability of the selected one-form carrier. Hypothesis \emph{(iii)}
says that no proper visible subcarrier still supports gauge-covariant,
exactness-compatible exchange response together with those same structural
properties. Hypothesis \emph{(iv)} says that any competing characteristic
observer datum satisfying \emph{(i)}--\emph{(iii)} is horizon-preserving
equivalent to the datum generated by the selected one-form carrier. Therefore
that datum satisfies the clauses of
Definition~\ref{def:least-motion-characteristic-observer} and is unique up to
horizon-preserving equivalence. Hence the selected one-form carrier is the
unique least-motion characteristic observer for the realized family.
\end{proof}

\begin{proof}[Proof of Theorem~\ref{thm:induced-gauge-compatible-law}]
\label{proof:thm:induced-gauge-compatible-law}
By Theorem~\ref{thm:least-motion-gauge-observer}, the visible law is carried by
the selected one-form carrier. Hypothesis~\emph{(i)} of that theorem and
Corollary~\ref{cor:intrinsic-package-yields-gauge-compatible-local-datum}
therefore place on that selected cut a unique gauge-compatible local internal
datum, while Corollary~\ref{cor:realized-local-gauge-covariance} gives local
gauge covariance of the selected transport.

Theorem~\ref{thm:temporal-realization-forbids-primitive-internal-creation}
permits only exchange-current source data on that cut. Thus the
compatibility condition \(\delta J=0\) is the one-form expression of exact
balanced internal exchange rather than an independent constitutive insertion.

Because the visible datum is a local gauge class \([A]\), the transported
curvature \(F:=dA\) is independent of the chosen representative. Exactness of
the underlying complex gives
\[
dF=d^2A=0.
\]
Definition~\ref{def:selected-visible-effective-law} now identifies the
selected endpoint closure class on the selected carrier as the exact
carrier-side closure class of the one selected visible law on that cut. The
closure statement therefore reduces to the unique
gauge-compatible response operator in that class; denote it by
\(\mathcal M_{\mathrm g}\). Its output is the selected exchange current plus the
remaining local gauge-invariant endpoint error, recorded by
\(\mathsf R_{\mathrm g}\). This gives the displayed gauge-compatible law.
\end{proof}

\begin{proof}[Proof of Theorem~\ref{thm:gauge-local-endpoint-collapse}]
\label{proof:thm:gauge-local-endpoint-collapse}
By Theorem~\ref{thm:induced-gauge-compatible-law}, there is already a unique
minimal gauge-compatible local response operator
\[
\mathcal M_{\mathrm g}:\Omega^1(M)/d\Omega^0(M)\to\Omega^1(M)
\]
on the selected cut, satisfying
\[
dF=0,
\qquad
\mathcal M_{\mathrm g}([A])=J+\mathsf R_{\mathrm g}.
\]

Because the visible datum is the local gauge class \([A]\), local gauge
covariance implies that \(\mathcal M_{\mathrm g}\) depends only on the
curvature \(F=dA\). Item~\emph{(iii)} says that the response family has
minimal local order \(2\) in \(A\), hence local order \(1\) in \(F\). Thus
there is a first-order local curvature-side operator
\[
\widetilde{\mathcal M}_{\mathrm g}:\Omega^2(M)\to\Omega^1(M)
\]
with
\[
\mathcal M_{\mathrm g}([A])=\widetilde{\mathcal M}_{\mathrm g}(F).
\]

Item~\emph{(iii)} and
Theorem~\ref{thm:refinement-coherent-core-realization} supply a constructive
local representative family for that already exact selected endpoint class, so
the reconstructed endpoint law has a unique exact core representative up to
local gauge and continuum equivalence. By
Theorem~\ref{thm:energy-orthogonal-cut-forces-canonical-duality}, applied to
the selected one-form carrier using item~\emph{(i)} and the already forced
uniqueness of the endpoint closure class from
Theorem~\ref{thm:induced-gauge-compatible-law}, the final carrier-level
representative is fixed internally by the selected-cut pairing.

Let \(\delta_{\chi}\) denote the adjoint of \(d\) with respect to that
canonical selected-cut pairing on the one-form carrier. Since the response has
already descended to a first-order local curvature operator and the
pairing-compatible strong representative is unique, the only admissible
pairing-compatible representative of that constructively forced endpoint class is
the curvature-adjoint response. Therefore
\[
\widetilde{\mathcal M}_{\mathrm g}(F)=\delta_{\chi}F,
\qquad\text{hence}\qquad
\mathcal M_{\mathrm g}([A])=\delta_{\chi}dA=\delta_{\chi}F.
\]
Substituting this identity into
Theorem~\ref{thm:induced-gauge-compatible-law} yields the displayed system.
The lossless statement is immediate.
\end{proof}

\begin{proof}[Proof of Corollary~\ref{cor:gauge-schur-remainder-identity}]
\label{proof:cor:gauge-schur-remainder-identity}
By Theorem~\ref{thm:canonical-selected-schur-realization}, the only hidden-side
visible forcing on the selected cut is the returned residual, and that forcing
is exactly the Schur correction \(\CCSigma{h}{z}\). By
Theorem~\ref{thm:gauge-local-endpoint-collapse}, the selected gauge endpoint
class has a unique canonical curvature-adjoint representative on the selected
carrier. Therefore the common Schur correction has a unique gauge-side image on
that representative; denote it by \(\Sigma_{h,\mathrm g}(z)\). By
Theorem~\ref{thm:returned-remainders-form-predictive-correction-surface}, the
gauge remainder is exactly that returned correction on the gauge lane, so
\(\mathsf R_{\mathrm g}=\Sigma_{h,\mathrm g}(z)\). Substituting this into
Theorem~\ref{thm:gauge-local-endpoint-collapse} gives the displayed identity.
\end{proof}

\begin{proof}[Proof of Theorem~\ref{thm:hidden-channel-gauge-current-and-branch-splitting}]
\label{proof:thm:hidden-channel-gauge-current-and-branch-splitting}
By item \emph{(v)} of
Definition~\ref{def:coupled-visible-shadow-gauge-mode-model},
\[
\bigl(\Omega_{\mathrm{g,vis}}(k)^2-\omega^2\bigr)a
=
j_0+\gamma_{\chi}b,
\]
\[
\bigl(\Omega_{\mathrm{g,hid}}(k)^2-\omega^2\bigr)b
=
\gamma_{\chi}a.
\]
If \(\Omega_{\mathrm{g,hid}}(k)^2\neq\omega^2\), the second identity gives
\[
b
=
\frac{\gamma_{\chi}}{\Omega_{\mathrm{g,hid}}(k)^2-\omega^2}\,a.
\]
Substituting this into the visible equation yields
\[
\bigl(\Omega_{\mathrm{g,vis}}(k)^2-\omega^2\bigr)a
=
j_0
+
\frac{\gamma_{\chi}^2}{\Omega_{\mathrm{g,hid}}(k)^2-\omega^2}\,a.
\]
This is exactly the displayed visible response law with
\[
\sigma_{\chi}^{\mathrm{g,hid}}(k,\omega)
=
\frac{\gamma_{\chi}^2}{\Omega_{\mathrm{g,hid}}(k)^2-\omega^2}.
\]
Multiplying by the mode one-form
\(a\,\varepsilon\,e^{i(nk\ell_{\chi}-\omega t)}\) gives the displayed
one-form remainder. If \(j_0=0\), then multiplying the visible equation by
\(\Omega_{\mathrm{g,hid}}(k)^2-\omega^2\) gives
\[
\bigl(\Omega_{\mathrm{g,vis}}(k)^2-\omega^2\bigr)
\bigl(\Omega_{\mathrm{g,hid}}(k)^2-\omega^2\bigr)
=
\gamma_{\chi}^2,
\]
equivalently
\[
(\omega^2-\Omega_{\mathrm{g,vis}}(k)^2)
(\omega^2-\Omega_{\mathrm{g,hid}}(k)^2)
=
\gamma_{\chi}^2.
\]
This is a quadratic equation in \(\omega^2\), and solving it gives the two
displayed branch functions.
\end{proof}

\begin{proof}[Proof of Corollary~\ref{cor:exact-off-resonant-gauge-pole-shifts}]
\label{proof:cor:exact-off-resonant-gauge-pole-shifts}
By Theorem~\ref{thm:hidden-channel-gauge-current-and-branch-splitting},
\[
\omega_{\pm}(k)^2
=
\frac{
\Omega_{\mathrm{g,vis}}(k)^2+\Omega_{\mathrm{g,hid}}(k)^2
\pm
\sqrt{
\bigl(\delta_{\chi}^{\mathrm g}(k)\bigr)^2+4\gamma_{\chi}^2
}
}{2}.
\]
Since
\(\Omega_{\mathrm{g,hid}}(k)^2=\Omega_{\mathrm{g,vis}}(k)^2+\delta_{\chi}^{\mathrm g}(k)\),
\[
\omega_-(k)^2
=
\Omega_{\mathrm{g,vis}}(k)^2
+
\frac{
\delta_{\chi}^{\mathrm g}(k)
-
\sqrt{
\bigl(\delta_{\chi}^{\mathrm g}(k)\bigr)^2+4\gamma_{\chi}^2
}
}{2}.
\]
Rationalizing the numerator gives the displayed formula for \(\omega_-(k)^2\),
and the formula for \(\omega_+(k)^2\) is obtained similarly. Because
\(\delta_{\chi}^{\mathrm g}(k)>0\), the visible-like shift is negative and the
hidden-like shift is positive.
\end{proof}

\begin{proof}[Proof of Corollary~\ref{cor:exact-avoided-crossing-coupled-gauge-branches}]
\label{proof:cor:exact-avoided-crossing-coupled-gauge-branches}
At a mode \(k_\ast\) with
\(\Omega_{\mathrm{g,vis}}(k_\ast)^2=\Omega_{\mathrm{g,hid}}(k_\ast)^2\), the
branch formula of
Theorem~\ref{thm:hidden-channel-gauge-current-and-branch-splitting} becomes
\[
\omega_{\pm}(k_\ast)^2
=
\frac{2\Omega_{\ast}^2\pm\sqrt{4\gamma_{\chi}^2}}{2}
=
\Omega_{\ast}^2\pm\gamma_{\chi}.
\]
Subtracting the two identities gives the exact gap.
\end{proof}

\begin{proof}[Proof of Corollary~\ref{cor:low-frequency-induced-exchange-susceptibility}]
\label{proof:cor:low-frequency-induced-exchange-susceptibility}
By Theorem~\ref{thm:hidden-channel-gauge-current-and-branch-splitting},
\[
\sigma_{\chi}^{\mathrm{g,hid}}(k,\omega)
=
\frac{\gamma_{\chi}^2}{
\Omega_{\mathrm{g,hid}}(k)^2-\omega^2
}.
\]
For small \(k\),
\[
\sin^2\!\left(\frac{k\ell_{\chi}}{2}\right)
=
\frac{k^2\ell_{\chi}^2}{4}+O(k^4),
\]
so
\[
\Omega_{\mathrm{g,hid}}(k)^2
=
\nu_{\chi}^2+\widetilde c_{\chi}^2k^2+O(k^4).
\]
Hence
\[
\Omega_{\mathrm{g,hid}}(k)^2-\omega^2
=
\nu_{\chi}^2+\widetilde c_{\chi}^2k^2-\omega^2+O(k^4).
\]
Writing
\[
x
:=
\frac{\widetilde c_{\chi}^2k^2-\omega^2+O(k^4)}{\nu_{\chi}^2},
\]
we have
\[
\sigma_{\chi}^{\mathrm{g,hid}}(k,\omega)
=
\frac{\gamma_{\chi}^2}{\nu_{\chi}^2}\frac{1}{1+x}.
\]
Expanding \((1+x)^{-1}=1-x+O(x^2)\) gives
\[
\sigma_{\chi}^{\mathrm{g,hid}}(k,\omega)
=
\frac{\gamma_{\chi}^2}{\nu_{\chi}^2}
-
\frac{\gamma_{\chi}^2\widetilde c_{\chi}^2}{\nu_{\chi}^4}k^2
+
\frac{\gamma_{\chi}^2}{\nu_{\chi}^4}\omega^2
+
O(k^4+k^2\omega^2+\omega^4),
\]
as claimed.
\end{proof}

% ── Chapter 8: flagship_geometry.tex (Geometric / Einstein) ──

\begin{proof}[Proof of Theorem~\ref{thm:least-motion-geometric-observer}]
\label{proof:thm:least-motion-geometric-observer}
Hypothesis \emph{(i)} gives faithfulness, closure-admissibility, and promotion
stability of the symmetric \(2\)-tensor carrier. Hypothesis \emph{(ii)} says
that no proper visible subcarrier still supports a faithful
divergence-compatible second-order response law at the observed scale.
Hypothesis \emph{(iii)} says that any competing characteristic observer datum
with those properties is horizon-preserving equivalent to the symmetric
\(2\)-tensor carrier. Therefore the datum generated by that carrier satisfies
the clauses of Definition~\ref{def:least-motion-characteristic-observer} and
is unique up to horizon-preserving equivalence. Hence the symmetric
\(2\)-tensor carrier is the unique least-motion characteristic observer for
the promoted source/transport family.
\end{proof}

\begin{proof}[Proof of Theorem~\ref{thm:weak-field-geometric-response-law}]
\label{proof:thm:weak-field-geometric-response-law}
The least-motion observer theorem places the observable law on the symmetric
tensor carrier. Interpreted through temporal realization and least-motion
selection, this carrier is the minimal faithful divergence-compatible visible
cut at the selected scale. By
Corollary~\ref{cor:realized-isotropic-symmetric-response-scalarization}, the
orthogonal-equivariance hypothesis forces the trace line and traceless carrier
to be the intrinsic scalar channels of the leading response. The
divergence-compatibility hypothesis therefore selects a unique leading operator
\(\mathcal E_{\mathrm{lin}}\) compatible with that two-channel decomposition,
while the residual scalar trace contribution is recorded by
\(\Lambda_{\mathrm{ind}}\). The source coupling to the preceding selected law
is, by
Theorem~\ref{thm:temporal-realization-forbids-primitive-internal-creation},
likewise an exchange coupling on the same selected channels,
hence is encoded by a single induced coefficient \(\kappa_{\mathrm{ind}}\).
All quadratic and higher response terms are absorbed into
\(\mathsf R_{\mathrm{grav}}\).
\end{proof}

\begin{proof}[Proof of Theorem~\ref{thm:characteristic-covariant-geometric-reduction}]
\label{proof:thm:characteristic-covariant-geometric-reduction}
By Theorem~\ref{thm:weak-field-geometric-response-law}, the weak-field law
already furnishes the induced coefficients \(\Lambda_{\mathrm{ind}}\) and
\(\kappa_{\mathrm{ind}}\) together with the leading divergence-compatible
linear response operator \(\mathcal E_{\mathrm{lin}}\). By hypothesis, the
nonlinear correction is carried by a rapidly decaying characteristic memory
integral. Corollary~\ref{cor:characteristic-boundary-reduction} removes its
bulk contribution, so after the stated slow-variation hypothesis only the
endpoint response remains at leading order.

The endpoint map acts on the symmetric tensor carrier. Under least-motion
selection on the temporally realized class, this is the selected covariant
visible cut, so all
unresolved geometric return must re-enter through its exact endpoint response.
By the no-preferred-direction clause and
Corollary~\ref{cor:realized-isotropic-symmetric-response-scalarization}, the
trace line and traceless carrier are the intrinsic scalar channels of that
endpoint response. Covariance and divergence compatibility therefore promote
the endpoint operator to a local covariant completion class extending the
weak-field operator \(\mathcal E_{\mathrm{lin}}+\Lambda_{\mathrm{ind}}I\). By
Definition~\ref{def:selected-visible-effective-law}, the selected cut already
carries the exact carrier-side endpoint closure class of the one selected
visible law, and its covariant representative has order \(2\) by the final
hypothesis. This yields a covariant
divergence-compatible response operator \(\mathcal G(g)\) of minimal local
order. Theorem~\ref{thm:temporal-realization-forbids-primitive-internal-creation}
identifies the divergence-compatible source field \(T\) as the visible
symmetric exchange tensor of the coupled sector. The same
scalar-channel decomposition therefore encodes that exchange coupling through
the induced coefficient
\(\kappa_{\mathrm{ind}}\). All higher endpoint corrections are local symmetric
tensors and are absorbed into \(\mathsf N_{\mathrm{grav}}(g,T)\). This yields
the displayed covariant law, and the divergence statement is exactly the
divergence-compatible endpoint identity together with \(\nabla g=0\).
\end{proof}

\begin{proof}[Proof of Theorem~\ref{thm:geometric-local-endpoint-collapse}]
\label{proof:thm:geometric-local-endpoint-collapse}
By Theorem~\ref{thm:characteristic-covariant-geometric-reduction}, there are
already induced coefficients \(\Lambda_{\mathrm{ind}},\kappa_{\mathrm{ind}}\),
a covariant divergence-compatible response operator \(\mathcal G(g)\) of
minimal local order, and a local symmetric remainder
\(\mathsf N_{\mathrm{grav}}(g,T)\) satisfying the displayed covariant law on
the selected cut.

Item~\emph{(ii)} says that the admissible covariant response family is
represented levelwise by a local symmetric-tensor differential family of
minimal local order \(2\), and item~\emph{(iii)} gives a refinement-coherent
exact-core realization of that family on its dense test domain. By
Theorem~\ref{thm:energy-orthogonal-cut-forces-canonical-duality}, applied to
the selected symmetric-tensor carrier using item~\emph{(i)} and the already
forced uniqueness of the geometric endpoint class from
Theorem~\ref{thm:characteristic-covariant-geometric-reduction}, the final
carrier-level representative is fixed internally by the selected-cut pairing.

By
Corollary~\ref{cor:realized-isotropic-symmetric-response-scalarization}, the
trace line and traceless carrier are the intrinsic scalar channels of that
response. Let \(\mathcal G_{\chi}(g)\) denote the unique pairing-compatible
divergence-compatible representative of that constructively forced endpoint
class on those two scalar channels, and let \(\operatorname{Div}_{\chi}\)
denote the corresponding selected-cut divergence operator on the symmetric
tensor carrier. Substituting this unique representative into
Theorem~\ref{thm:characteristic-covariant-geometric-reduction} yields
\[
\mathcal G_{\chi}(g)+\Lambda_{\mathrm{ind}}g
=
\kappa_{\mathrm{ind}}T+\mathsf N_{\mathrm{grav}}(g,T),
\qquad
\operatorname{Div}_{\chi}\bigl(\mathcal G_{\chi}(g)+\Lambda_{\mathrm{ind}}g\bigr)=0.
\]
This is exactly the displayed intrinsic selected-cut geometric law. The
lossless statement is immediate.
\end{proof}

\begin{proof}[Proof of Corollary~\ref{cor:geometric-schur-remainder-identity}]
\label{proof:cor:geometric-schur-remainder-identity}
By Theorem~\ref{thm:canonical-selected-schur-realization}, the selected cut
carries one common returned residual field and that returned residual is
exactly the Schur correction \(\CCSigma{h}{z}\). By
Theorem~\ref{thm:geometric-local-endpoint-collapse}, the selected geometric
endpoint class has a unique canonical trace-adjusted curvature representative on the selected
carrier. Therefore the common Schur correction has a unique geometric-side
image on that representative; denote it by \(\Sigma_{h,\mathrm{grav}}(z)\). By
Theorem~\ref{thm:returned-remainders-form-predictive-correction-surface}, the
geometric remainder is exactly that returned correction on the geometric lane,
so \(\mathsf N_{\mathrm{grav}}(g,T)=\Sigma_{h,\mathrm{grav}}(z)\). Substituting
this into Theorem~\ref{thm:geometric-local-endpoint-collapse} yields the
displayed equation.
\end{proof}

\begin{proof}[Proof of Theorem~\ref{thm:hidden-channel-geometric-remainder-and-branch-splitting}]
\label{proof:thm:hidden-channel-geometric-remainder-and-branch-splitting}
By item \emph{(v)} of
Definition~\ref{def:coupled-visible-shadow-geometric-mode-model},
\[
\bigl(\Omega_{\mathrm{grav,vis}}(k)^2-\omega^2\bigr)u
=
\tau_0+\gamma_{\chi}v,
\]
\[
\bigl(\Omega_{\mathrm{grav,hid}}(k)^2-\omega^2\bigr)v
=
\gamma_{\chi}u.
\]
If \(\Omega_{\mathrm{grav,hid}}(k)^2\neq\omega^2\), the second identity gives
\[
v
=
\frac{\gamma_{\chi}}{\Omega_{\mathrm{grav,hid}}(k)^2-\omega^2}\,u.
\]
Substituting this into the visible equation yields
\[
\bigl(\Omega_{\mathrm{grav,vis}}(k)^2-\omega^2\bigr)u
=
\tau_0
+
\frac{\gamma_{\chi}^2}{\Omega_{\mathrm{grav,hid}}(k)^2-\omega^2}\,u.
\]
This is exactly the displayed visible response law with
\[
\sigma_{\chi}^{\mathrm{grav,hid}}(k,\omega)
=
\frac{\gamma_{\chi}^2}{\Omega_{\mathrm{grav,hid}}(k)^2-\omega^2}.
\]
Multiplying by the mode tensor
\(u\,E\,e^{i(nk\ell_{\chi}-\omega t)}\) gives the displayed symmetric-tensor
remainder. If \(\tau_0=0\), then multiplying the visible equation by
\(\Omega_{\mathrm{grav,hid}}(k)^2-\omega^2\) gives
\[
\bigl(\Omega_{\mathrm{grav,vis}}(k)^2-\omega^2\bigr)
\bigl(\Omega_{\mathrm{grav,hid}}(k)^2-\omega^2\bigr)
=
\gamma_{\chi}^2,
\]
equivalently
\[
(\omega^2-\Omega_{\mathrm{grav,vis}}(k)^2)
(\omega^2-\Omega_{\mathrm{grav,hid}}(k)^2)
=
\gamma_{\chi}^2.
\]
This is a quadratic equation in \(\omega^2\), and solving it gives the two
displayed branch functions.
\end{proof}

\begin{proof}[Proof of Corollary~\ref{cor:exact-off-resonant-geometric-pole-shifts}]
\label{proof:cor:exact-off-resonant-geometric-pole-shifts}
By Theorem~\ref{thm:hidden-channel-geometric-remainder-and-branch-splitting},
\[
\omega_{\pm}(k)^2
=
\frac{
\Omega_{\mathrm{grav,vis}}(k)^2+\Omega_{\mathrm{grav,hid}}(k)^2
\pm
\sqrt{
\bigl(\delta_{\chi}^{\mathrm{grav}}(k)\bigr)^2+4\gamma_{\chi}^2
}
}{2}.
\]
Since
\(\Omega_{\mathrm{grav,hid}}(k)^2=\Omega_{\mathrm{grav,vis}}(k)^2+\delta_{\chi}^{\mathrm{grav}}(k)\),
\[
\omega_-(k)^2
=
\Omega_{\mathrm{grav,vis}}(k)^2
+
\frac{
\delta_{\chi}^{\mathrm{grav}}(k)
-
\sqrt{
\bigl(\delta_{\chi}^{\mathrm{grav}}(k)\bigr)^2+4\gamma_{\chi}^2
}
}{2}.
\]
Rationalizing the numerator gives the displayed formula for \(\omega_-(k)^2\),
and the formula for \(\omega_+(k)^2\) is obtained similarly. Because
\(\delta_{\chi}^{\mathrm{grav}}(k)>0\), the visible-like shift is negative and
the hidden-like shift is positive.
\end{proof}

\begin{proof}[Proof of Corollary~\ref{cor:exact-avoided-crossing-coupled-geometric-branches}]
\label{proof:cor:exact-avoided-crossing-coupled-geometric-branches}
At a mode \(k_\ast\) with
\(\Omega_{\mathrm{grav,vis}}(k_\ast)^2=\Omega_{\mathrm{grav,hid}}(k_\ast)^2\),
the branch formula of
Theorem~\ref{thm:hidden-channel-geometric-remainder-and-branch-splitting}
becomes
\[
\omega_{\pm}(k_\ast)^2
=
\frac{2\Omega_{\ast}^2\pm\sqrt{4\gamma_{\chi}^2}}{2}
=
\Omega_{\ast}^2\pm\gamma_{\chi}.
\]
Subtracting the two identities gives the exact gap.
\end{proof}

\begin{proof}[Proof of Corollary~\ref{cor:low-frequency-induced-geometric-susceptibility}]
\label{proof:cor:low-frequency-induced-geometric-susceptibility}
By Theorem~\ref{thm:hidden-channel-geometric-remainder-and-branch-splitting},
\[
\sigma_{\chi}^{\mathrm{grav,hid}}(k,\omega)
=
\frac{\gamma_{\chi}^2}{
\Omega_{\mathrm{grav,hid}}(k)^2-\omega^2
}.
\]
For small \(k\),
\[
\sin^2\!\left(\frac{k\ell_{\chi}}{2}\right)
=
\frac{k^2\ell_{\chi}^2}{4}+O(k^4),
\]
so
\[
\Omega_{\mathrm{grav,hid}}(k)^2
=
\nu_{\chi}^2+\widetilde c_{\chi}^2k^2+O(k^4).
\]
Hence
\[
\Omega_{\mathrm{grav,hid}}(k)^2-\omega^2
=
\nu_{\chi}^2+\widetilde c_{\chi}^2k^2-\omega^2+O(k^4).
\]
Writing
\[
x
:=
\frac{\widetilde c_{\chi}^2k^2-\omega^2+O(k^4)}{\nu_{\chi}^2},
\]
we have
\[
\sigma_{\chi}^{\mathrm{grav,hid}}(k,\omega)
=
\frac{\gamma_{\chi}^2}{\nu_{\chi}^2}\frac{1}{1+x}.
\]
Expanding \((1+x)^{-1}=1-x+O(x^2)\) gives
\[
\sigma_{\chi}^{\mathrm{grav,hid}}(k,\omega)
=
\frac{\gamma_{\chi}^2}{\nu_{\chi}^2}
-
\frac{\gamma_{\chi}^2\widetilde c_{\chi}^2}{\nu_{\chi}^4}k^2
+
\frac{\gamma_{\chi}^2}{\nu_{\chi}^4}\omega^2
+
O(k^4+k^2\omega^2+\omega^4),
\]
as claimed.
\end{proof}
